% ===== 原文 pp.47--50（tr-03 第 3 段；本块到此结束，p.51 的 (5.24) 起归 tr-04）=====
% [原文 p.47]

\subsection{超对称局域化}

把超对称局域化用在拓扑扭曲场论的标量超荷上，已有几十年历史。它获得新生始于 Pestun 在 2007 年 [10] 对 4d $\cN=2$ 理论的球面配分函数以及 Wilson 圈期望值的计算。此后十年间，这一技术被成功应用于许多维数（从 1d 到 7d，甚至连续维数）与几何（例如球面 $S^d$、乘积 $S^{d-1}\times S^1$、半球以及其他带边界的空间），2016 年的综述文集 [125] 对此有总结。除了用于理解超对称理论与黑洞态计数之外，所得表达式往往还是本身即有价值的复杂特殊函数。我们在这里介绍这一技术，并在 5.4 小节给出一个变体，用以解释种种因子化性质。

\para{超对称局域化的设置。} 我们的目标是计算一个路径积分
\begin{equation}
  \langle\mathcal{O}\rangle=\int[\mathrm{D}\phi]\,\ee^{-S}\mathcal{O}
  \tag{5.10}
\end{equation}
它在一部分超荷 $Q$ 下不变（我们把所有场统记为 $\phi$）。这意味着作用量与路径积分测度是 $Q$ 不变的（$QS=0$ 且 $\int[\mathrm{D}\phi]\,Q(\text{任何东西})=0$），可观测量也是（$Q\mathcal{O}=0$）。粗略地说，被积函数在构型空间中沿 $Q$ 的轨道不变，于是每条非平凡轨道上的积分只是常数的 Grassmann 积分，因而为零。

\begin{keybox}{超对称局域化的配方}
\begin{itemize}[leftmargin=1.4em,itemsep=1pt,topsep=1pt]
\item 在 $M$ 上选择一个超引力背景，其中至少含一个广义共形 Killing 旋量 $\xi$，使 $M$ 上的理论至少有一个超对称 $Q$。要把它在壳外（off-shell）实现。
\item 选择一个费米型泛函 $V$，它 $Q^2$ 不变，且在路径积分围道上满足 $(QV)_{\text{bosonic}}\ge0$。
\item 求 $(QV)_{\text{bosonic}}$ 的零点，它们构成积分轨迹，往往是有限维的。
\item 在这些零点附近把 $(QV)_{\text{bosonic}}$ 展开到二次阶，算出高斯积分（单圈行列式）。
\end{itemize}
\end{keybox}

直观上，这应当把积分约化到 $Q$ 不变的场构型上。

超对称局域化以路径积分为基础，因此需要超对称 $Q$ 在壳外（而不只是在壳上）实现。对超对称性很高的理论（例如 4d $\cN=2$），要在壳上实现全部超对称需要无穷多个辅助场，但通常只需有限多个辅助场，就能把单个超荷在壳外实现。

超对称局域化的核心思想是：变形作用量 $S$，使 $\langle\mathcal{O}\rangle$ 不受影响、同时压制大多数构型的贡献，从而把路径积分约化到更小的构型空间上。通常我们会设法让这个更小的空间是有限维的，以便得到一个良定义的有限维积分；但得到一个更低维的场论有时也很有意思。

具体地说，需要某个场的泛函具备三条性质：它是 $Q$-恰当的（即形如 $QV$）、$Q$-闭的（即 $V$ 在玻色对称性 $Q^2$ 下不变），并且在我们所考虑的路径积分围道上玻色部分非负。一个典型选择大致是对理论中所有费米子（统记为 $\psi$）求和，形如
\begin{equation}
  V=\sum_{\text{所有旋量 }\psi}\psi\,\overline{Q\psi}
  \qquad\Longrightarrow\qquad
  (QV)_{\text{bosonic}}=\sum_{\psi}|Q\psi|^2
  \tag{5.11}
\end{equation}
其中共轭 $\overline{Q\psi}$ 的定义要选得合适，以确保 $(Q\psi)\overline{Q\psi}$ 沿路径积分围道为正。

% [Tutor Note: p.47 段末"沿 $Q$ 轨道被积函数不变 ⇒ 积分为零"可具体化：设沿某条非平凡轨道取费米坐标 $\eta$，被积函数与 $\eta$ 无关，则该方向的积分是 $\int\dd\eta\cdot c=0$（Berezin 积分中常数项积分为零）。所以只有 $Q$ 的不动点（轨道退化为一点）才有非零贡献------这正是 (5.13) 把积分压到 $QV=0$ 上的原因。]

% [原文 p.48]

\para{鞍点计算。} 一旦选定了这样一项，就用 $tQV$（$t\in[0,+\infty)$）变形作用量，并注意到可观测量不受影响，因为
\begin{equation}
\begin{aligned}
  \langle\mathcal{O}\rangle_t&=\int[\mathrm{D}\phi]\,\ee^{-S-tQV}\mathcal{O}\\
  \Longrightarrow\quad
  \partial_t\langle\mathcal{O}\rangle_t
  &=-\int[\mathrm{D}\phi]\,\ee^{-S-tQV}\mathcal{O}\,QV
   =-\int[\mathrm{D}\phi]\,Q\!\left(\ee^{-S-tQV}\mathcal{O}\,V\right)=0.
\end{aligned}
  \tag{5.12}
\end{equation}
这里我们用到了 $Q\mathcal{O}=0=Q(S+tQV)$，从而把被积函数写成某个东西的 $Q$。

可观测量与 $t$ 无关，因此可以取极限 $t\to\infty$；在这个极限下鞍点近似变成精确的。此外，任何 $(QV)_{\text{bosonic}}>0$ 的鞍点都被 $\ee^{-tQV}$ 无穷地压制。既然我们假设了 $QV\ge0$，在这极限下就只剩下 $QV=0$ 的场构型上的积分，以及其周围二次涨落的高斯积分：
\begin{equation}
  \langle\mathcal{O}\rangle=\langle\mathcal{O}\rangle_{t=0}
  =\lim_{t\to\infty}\langle\mathcal{O}\rangle_t
  =\int_{QV=0}[\mathrm{D}\phi]\,\ee^{-S[\phi]}Z_{\text{one-loop}}[\phi]\,\mathcal{O}[\phi].
  \tag{5.13}
\end{equation}
这里 $\int_{QV=0}$ 只是示意的写法；若 $QV$ 的零点空间不连通，它也可能包含离散求和。这里 $Z_{\text{one-loop}}[\phi]$ 是在 $QV$ 的一个零点构型 $\phi$ 附近对 $\exp(-tQV)$ 作高斯积分的结果。

计算单圈行列式 $Z_{\text{one-loop}}$ 往往是最技术性的一步。它由 $QV$ 中的二次项 $(QV)_2$ 决定，更准确地说，是 $(QV)_2$ 的费米部分与玻色部分的行列式之比。在足够简单的几何中，可以把所有模式列出来、显式对角化这些算子，并用例如 zeta 函数正则化来正规化对模式的无穷乘积。由于超荷 $Q$ 把大多数玻色模式与费米模式配成对，这一计算会出现大量相消。粗略地说，只有被 $Q$ 湮灭的模式才有贡献；这一观察有助于扩展那些可以用“初等”办法算出单圈行列式的情形。最后，用强有力的指标定理还可以把后一计算大大简化。细节请参看综述文集 [125]。

\begin{exercise}{5.3}
\begin{enumerate}[label=\arabic*.,leftmargin=2em,itemsep=1pt,topsep=1pt]
\item 回忆一下：对对称 $n\times n$ 矩阵 $\Delta$，有高斯积分
  $\int_{\R^n}\dd^n x\,\exp(-x^T\Delta x)=\sqrt{\det(\pi/\Delta)}$。
\item 再考虑 $n$ 对 Grassmann 变量 $\theta_1,\bar\theta_1,\dots,\theta_n,\bar\theta_n$ 与一个 $n\times n$ 矩阵 $D$，计算 Berezin 积分
  $\int\dd^n\theta\,\dd^n\bar\theta\,\exp(-\bar\theta\cdot D\cdot\theta)=\pm\det D$。
\end{enumerate}
在实际的 QFT 应用中，玻色算子 $\Delta$ 大致就是 Laplacian，费米算子 $D$ 大致就是 Dirac 算子，故 $\Delta\sim D^2$，它们的单圈行列式大部分相消。
\end{exercise}

% [原文 p.49]

\para{椭球上的鞍点。} 现在把局域化应用到压扁球面 $S^4_b$ 上的 4d $\cN=2$ 理论。我们取标准的变形项 (5.11)，其中求和跑遍夸克 $\psi$（超多重态旋量）与规范微子 $\lambda$（矢量多重态旋量）。所得的 $QV$ 与这些多重态的 4d $\cN=2$ 作用量相当接近，我们只提及确定空间 $(QV)_{\text{bosonic}}=0$ 所需的部分。

超对称局域化依赖于存在一个作为壳外对称性的超荷 $Q$。这需要给 4d $\cN=2$ 理论添加一些辅助场 $K$。眼下我们先专注于球面 [10]，只在最终表达式里恢复压扁 [11]。

对超多重态，我们有
\begin{equation}
  (QV_{\text{hyper}})_{\text{bosonic}}
  =|Dq|^2+|D\tilde{q}|^2+\dots+\frac{R}{6}|q|^2+\frac{R}{6}|\tilde{q}|^2+|K_q|^2.
  \tag{5.14}
\end{equation}
这里 $R$ 是 Ricci 标量，取正值：这一项来自把 $|Dq|^2$ 从平直空间共形映射到球面。式中的“$\dots$”也是一些平方之和，因此零点轨迹把整个超多重态都置零：
\begin{equation}
  q=\tilde{q}=K_q=0,
  \tag{5.15}
\end{equation}
费米子亦然，因为它们是 Grassmann 变量。

对矢量多重态，我们有类似的项，只是把 $q,\tilde{q}$ 换成了矢量多重态标量 $\phi$；但此外还有形如 $|F_{\mu\nu}|^2$ 的项，以及来自超引力背景的项。最终，（变形项的玻色部分）可以整理成如下平方和：
\begin{equation}
\begin{aligned}
  (QV)_{\text{vector,bosonic}}
  &=\frac{r-x^0}{2r}\big(F^-_{\mu\nu}+w^-_{\mu\nu}\Re\phi\big)^2
   +\frac{r+x^0}{2r}\big(F^+_{\mu\nu}+w^+_{\mu\nu}\Re\phi\big)^2\\
  &\quad+|D\phi|^2+[\phi,\phi^\dagger]^2+|K_{\phi,i}+w_i\Im\phi|^2.
\end{aligned}
  \tag{5.16}
\end{equation}
这里 $F^-$ 与 $F^+$ 是规范场强的（反）自对偶部分，$K_{\phi,i}$（$i=1,2,3$）是辅助场（$\fsu(2)_{\mathrm{R}}$ 的三重态），而 $w^\pm_{\mu\nu}$ 与 $w_i$ 由超引力背景决定。

让我们求 (5.16) 的零点。每一项都必须为零，特别地 $\phi$ 协变常定（$D\phi=0$）。在极点 $x^0=\pm r$ 之外，我们有 $F^\pm_{\mu\nu}=-w^\pm_{\mu\nu}\Re\phi$，于是 $F_{\mu\nu}=-w_{\mu\nu}\Re\phi$，其中 $w=w^++w^-$。接着 Bianchi 恒等式给出
\begin{equation}
  0=D^\mu F_{\mu\nu}=D^\mu(-w_{\mu\nu}\Re\phi)
  =-(\partial^\mu w_{\mu\nu})\Re\phi-w_{\mu\nu}\Re(D^\mu\phi)
  \tag{5.17}
\end{equation}
而末项因 $D^\mu\phi=0$ 而消失。在我们这里的特定超引力背景中 $\partial^\mu w_{\mu\nu}\ne0$，于是可知 $\Re\phi=0$，从而 $F_{\mu\nu}=0$。因此在合适的规范下规范场为零，并且
\begin{equation}
  A_\mu=0,\qquad \phi=a,\qquad K_{\phi,i}=-w_i a
  \qquad\text{（极点之外）},
  \tag{5.18}
\end{equation}
其中 $a$ 是常数。

% [原文 p.50]

另一方面，在极点 $x^0=\pm r$ 处，两个方程 $F^\pm_{\mu\nu}=-w^\pm_{\mu\nu}\Re\phi$ 中只有一个成立，而 $F_{\mu\nu}$ 的另一部分不受约束。这提示我们在极点处加入点状的瞬子构型：
\begin{equation}
  \text{瞬子 }(F^+=0)\text{ 在 }x^0=r;\qquad
  \text{反瞬子 }(F^-=0)\text{ 在 }x^0=-r.
  \tag{5.19}
\end{equation}

让我们专注于北极 $x^0=r$。瞬子构型对物质不敏感，因此瞬子模空间 $\cM_{\text{inst}}$ 可按单规范群因子分解为各规范群的瞬子模空间之积。后者又分裂为无穷多个连通分支之并，以瞬子数 $k=\#\int\Tr(F\wedge F)\in\Z_{\ge0}$ 标记（$\#$ 是某个可算的常数），每个规范群各有一个瞬子数。合起来，记规范群为 $\prod_I G_I$，
\begin{equation}
  \cM_{\text{inst}}=\prod_I\left(\bigsqcup_{k_I\ge0}\cM_{G_I,k_I}\right).
  \tag{5.20}
\end{equation}

\para{Coulomb 分支参数与质量的实性。} 条件 $\Re a=0$ 的意思很简单：$a\in\fg$ 而非其复化。我们把它规范固定到 Cartan 子代数 $\fh$（模 Weyl 群）中。于是实性条件就是
\begin{equation}
  \langle w,a\rangle\in\ii\R,\qquad\text{对每个权 }w.
  \tag{5.21}
\end{equation}
例如对 $\fg=\fsu(N)$，$a$ 是反 Hermite 的。

这个实性条件在这里是从研究变形项的鞍点得到的，但同一条件也可以从要求构型 $Q$ 不变推出。如前所述，超多重态质量 $m$ 无非是矢量多重态标量的常数背景值。因此，为使质量保持超对称局域化所用的超荷 $Q$，(5.21) 这个实性条件对质量也同样必需。换句话说，$S^4_b$ 上的弯曲空间超对称要求质量是虚的。

% [Tutor Note: p.50 式 (5.21) 的 $\langle w,a\rangle\in\ii\R$ 与"$a$ 在 $\fg$ 中"等价：取 $\fg=\fsu(N)$、$a=\ii\,\diag(a_1,\dots,a_N)$ 且 $\sum a_i=0$，权 $w$ 取 $e_i-e_j$ 的形式，则 $\langle w,a\rangle=\ii(a_i-a_j)\in\ii\R$ 自动成立；反之若某个 $a_k$ 有实部，则取相应的权即可破坏该条件。这就是「$a$ 反 Hermite」的具体含义，也是 (5.22) 中 $\Tr(ra)^2$ 为负实数（故 $-\Re(2\pi\ii\tau)\Tr(ra)^2$ 为实数作用量）的原因。]

\para{圆球面与压扁球面上的结果。} 由 (5.15)、(5.18)、(5.19) 定义的鞍点构型，其刻画方式是：对每个规范群选定一个虚的 Coulomb 分支参数 $a$，以及在极点处的点状（反）瞬子构型。对任何这样的鞍点，我们计算经典作用量
\begin{equation}
  S_{\mathrm{cl}}=-\Re(2\pi\ii\tau)\Tr(ra)^2+2\pi\ii\tau n-2\pi\ii\bar\tau\bar n,
  \qquad
  \tau=\frac{\theta}{2\pi}+\frac{4\pi\ii}{g^2},
  \tag{5.22}
\end{equation}
其中（压扁）球面 (5.8) 的半径 $r$ 使第一项无量纲。

% [Tutor Note: p.50 式 (5.22) 的写法 $\Tr(ra)^2$ 应读作 $\Tr\big((ra)^2\big)=r^2\Tr a^2$，而不是 $(\Tr ra)^2$------后者对 $\fsu(N)$ 恒为零（$a$ 无迹），且与"$r$ 使这一项无量纲"的说法不符（$a$、$r$ 的量纲分别是 $1/L$、$L$，只有 $ra$ 无量纲才需 $r^2\Tr a^2$）。原文把平方写在了括号外，是常见简写。另注：$2\pi\ii\tau n$ 与 $-2\pi\ii\bar\tau\bar n$ 中 $\bar\tau$、$\bar n$ 上的两条横线在原印本中几乎连成一笔，按物理（瞬子的复共轭贡献）应读作 $\bar\tau\bar n$。]

单圈行列式（高斯积分）刻划鞍点附近二次涨落的效应，其计算请参看 [126]。对零（反）瞬子鞍点，得到
\begin{equation}
  Z_{\text{one-loop}}=Z^{\text{vector}}_{\text{one-loop}}Z^{\text{hyper}}_{\text{one-loop}}
  \qquad\text{当 }k=\bar{k}=0.
  \tag{5.23}
\end{equation}

% [Tutor Note: p.50 式 (5.23) 只管 $k=\bar k=0$（零瞬子）的情形；$k\ne0$ 时单圈行列式还要乘上瞬子模空间的贡献，正是 5.3 小节 Nekrasov 配分函数要处理的对象。核对办法：把 (5.23) 与 (5.20) 对照，$k$ 与 $\bar k$ 分别是标记在 $x^0=r$ 与 $x^0=-r$ 两端的瞬子数，二者独立；这也解释了 (5.22) 中为何出现 $n$ 与 $\bar n$ 两项。]
